\begin{abstract}
Some studies on classifying functional (FR) and non-functional (NFR) requirements report accuracies of 98\% to 99\%. We examined three distinct evaluation concerns: duplicate rows shared across random folds, oversampling before splitting, and generalisation to unseen projects. Only the first two are leakage mechanisms; project-grouped evaluation measures performance on a different deployment setting. We measured their effects on PROMISE\_exp (969 requirements, 47 projects), with a replication on the Dalpiaz et al.\ data. We evaluated seven classifiers, including a dummy baseline, hard voting and stacking, under stratified, duplicate-grouped and project-grouped 5-fold CV repeated with 5 seeds. Model comparisons on the same splits used the corrected repeated $k$-fold $t$-test with Holm adjustment. Differences between protocols are reported descriptively. The protocol was locked before any run. Oversampling before splitting raised macro-F1 on the 12-class task by 0.39 to 0.49, to between 0.95 and 0.99. On the binary task the increase was 0.02 to 0.03. Testing on unseen projects lowered macro-F1 by 0.13 to 0.19 and changed the model ranking. Grouping the few duplicates changed macro-F1 by less than 0.01, and the ensembles beat only the random forest. We recommend oversampling inside the training fold, reporting duplicate handling and project-level results, and testing on the same folds before ranking models.
\end{abstract}

\begin{IEEEkeywords}
requirements classification, functional and non-functional requirements, data leakage, evaluation methodology, oversampling, ensemble learning
\end{IEEEkeywords}
